\documentclass[12pt]{article}
\usepackage[margin=1in]{geometry}
\usepackage[T1]{fontenc}
\usepackage[utf8]{inputenc}
\usepackage{lmodern}
\usepackage{setspace}
\usepackage{hyperref}
\usepackage{cite}
\usepackage{amsmath}
\usepackage{booktabs}
\usepackage{graphicx}
\usepackage{microtype}
\usepackage{amssymb}
\setstretch{1.05}
\hypersetup{colorlinks=true,linkcolor=blue,citecolor=blue,urlcolor=blue}

\title{Interpreting Transformer Representations on Rubik's Cube Tasks\\
\large Progress Report}
\author{Cole McGuire}
\date{}

\begin{document}
\maketitle

% ============================================================
\begin{abstract}
% ============================================================
We use the $2{\times}2{\times}2$ Rubik's cube as a controlled, fully-enumerable
domain for mechanistic interpretability.  A small transformer
($d_\text{model}{=}128$, 4 layers, ${\sim}200\text{k}$ parameters) is trained to
classify the optimal half-turn-metric distance to solved (0--11), achieving
78.5\% validation accuracy against an 8.3\% random baseline.  We then apply
four interpretability methods to understand what the model has learned.  Linear
probing finds that face-level solved status is 98\% decodable at \emph{every}
layer including the raw embedding, and that optimal distance is encoded almost
completely in the embedding itself (MAE drops from 0.87 to 0.40 in the first
block).  Activation patching reveals a dissociation: the probe-readable features
are not causally active---ablating their directions produces zero change in
model output---whereas swapping the full residual stream between distance classes
achieves a 100\% prediction flip at every layer.  Tuned-lens analysis reconciles
these findings: logit-lens accuracy grows from 20\% to 78\% across blocks,
demonstrating that each block is contributing meaningful computation, but in a
rotated basis that the final unembedding layer alone cannot read.  Sparse
autoencoders confirm that face and distance features are represented
monosemantically ($r \approx 0.81$--$0.89$), while corner-orientation features
are weaker and concentrated at the embedding layer.  Finally, motivated by
instructor feedback, we evaluate how well the cube state can be conveyed to a
pre-trained language model through text.  Of four representations tested, a
new \emph{piece-identity} encoding---which explicitly names each corner cubie by
its solved-state colors and states its current position and orientation---is the
only state-based representation to consistently solve depth-5 scrambles, though
depth-7 remains a hard wall for all representations including chain-of-thought
prompting.
\end{abstract}

% ============================================================
\section{Introduction}
% ============================================================

Mechanistic interpretability asks not only whether a neural network can solve a
task, but \emph{how}: which features are encoded in intermediate representations,
which are causally relevant to the output, and how computation flows through
layers.  Most such work has been done on large language models, where the sheer
scale makes it difficult to achieve a complete account.  We instead choose a
domain where the full state space is enumerable, optimal solutions are known for
every state, and the relevant geometric structure is explicit: the $2{\times}2{\times}2$
Rubik's cube.

Prior work has explored the cube primarily from a performance perspective.
Agostinelli et al.\ train a deep reinforcement learning agent that solves all
test configurations \cite{agostinelli2019deepcubea}.  Takano demonstrates that
self-supervised training on reverse scrambles is sufficient for near-optimal
solving \cite{takano2021selfsupervision}.  Chasmai proposed CubeTR, a transformer trained via sequence modelling on the
cube, though this work was subsequently withdrawn for lacking experimental
validation \cite{chasmai2021cubetr}.  Gupta et al.\ provide a more rigorous
transformer-based study, using the $2{\times}2{\times}2$ cube as a controlled
domain to investigate how explicit world-model pretraining objectives affect
representations and downstream reinforcement learning, combining linear probing
and causal interventions with GRPO post-training
\cite{gupta2024worldmodels}.  Our work is complementary: where Gupta et al.\
compare training objectives and study the effect on solving performance, we fix
the training objective and apply a broader suite of interpretability tools to
characterize what a single trained model represents.

For interpretability methodology we follow Alain and Bengio's framework of
linear classifier probes \cite{alain2016probes}, Belinkov's critical survey of
their promises and limitations \cite{belinkov2022probing}, Belrose et al.'s
tuned-lens technique for reading intermediate predictions \cite{belrose2023tunedlens},
Meng et al.'s causal activation-patching methodology
\cite{meng2022locatingfactual}, Bricken et al.'s sparse autoencoder approach
to monosemantic feature decomposition \cite{bricken2023monosemantic}, and
Nanda et al.'s analysis of grokking via mechanistic interpretability
\cite{nanda2023grokking}.  We are also motivated by Variengien et al.'s
theoretical result that belief-state geometry is linearly represented in
transformer residual streams \cite{variengien2024beliefstate}---the cube
provides a clean, discrete instance of this claim.

The core question this project addresses is: \emph{what internal representations
does a small transformer learn when trained on a combinatorial task with known
optimal structure, and do those representations match the task's known
geometric features?}  A secondary question, which emerged from instructor
feedback mid-project, is: \emph{what text representation of cube state is
necessary and sufficient for a pre-trained language model to reason about the
cube?}

% ============================================================
\section{Related Work}
% ============================================================

\paragraph{Cube-solving with neural networks.}
DeepCubeA \cite{agostinelli2019deepcubea} uses a value network trained by
BFS-bootstrapping; Takano \cite{takano2021selfsupervision} shows reverse-scramble
self-supervision yields state-of-the-art performance.  Chasmai proposed CubeTR, a transformer trained via sequence modelling, though
this work was withdrawn for lacking experimental validation
\cite{chasmai2021cubetr}.  Most directly related is Gupta et al.\
\cite{gupta2024worldmodels}, who train transformers on the $2{\times}2{\times}2$
cube with explicit world-model objectives and apply linear probes and causal
interventions---methodology closely parallel to ours.  Our work differs in focus:
we analyze a single classification model with a richer interpretability suite
(tuned lens, SAE, text representation testing) rather than comparing training
objectives.

\paragraph{Linear probing.}
Alain and Bengio \cite{alain2016probes} propose training linear classifiers on
frozen intermediate representations to measure what information each layer
encodes.  Belinkov \cite{belinkov2022probing} surveys methodological pitfalls,
including the risk of confusing \emph{encoding} with \emph{use}.  We address this
risk directly via activation patching (Section \ref{sec:patch}).

\paragraph{Tuned lens.}
The logit lens \cite{nostalgebraist2020logitlens} projects intermediate residual
stream states directly through the final unembedding.  Belrose et al.\
\cite{belrose2023tunedlens} improve this with a per-layer affine probe (tuned
lens), showing it is more reliable and less biased.  We apply both to our cube
model.

\paragraph{Activation patching.}
Meng et al.\ \cite{meng2022locatingfactual} introduce causal mediation analysis
for transformers: swapping activations from a clean to a corrupted run identifies
which components are causally responsible for a prediction.  We adapt this to a
concept-direction ablation: we project probe directions out of the residual stream
and measure whether model output changes.

\paragraph{Sparse autoencoders.}
Bricken et al.\ \cite{bricken2023monosemantic} train overcomplete autoencoders
on transformer residual streams to find monosemantic features that do not emerge
from direct examination of weight matrices.  We apply this technique to each
layer of our cube model.

\paragraph{Belief state geometry.}
Variengien et al.\ \cite{variengien2024beliefstate} prove that optimal
next-token predictors represent belief states linearly in their residual streams,
even when those belief states have fractal geometry.  The cube's state space is a
discrete dynamical system, and a model processing a cube state token should
represent a point mass over the 88M-state space.  Our probing results can
be read as an empirical test of their theoretical prediction.

% ============================================================
\section{Methods}
% ============================================================

\subsection{Cube Simulator and State Encoding}

The $2{\times}2{\times}2$ cube has 24 stickers arranged on 6 faces.  We assign each
face a canonical color (white top, yellow bottom, green front, blue back, orange
left, red right) and index stickers 0--23 in reading order per face (U=0--3,
D=4--7, F=8--11, B=12--15, L=16--19, R=20--23).  The state is a
$(24,)$ int8 array of color indices.  We encode this as a $(144,)$ float32
one-hot vector: $\mathbf{x} = \text{flatten}(\mathbf{I}_6[\mathbf{s}])$, where
$\mathbf{s}$ is the state array and $\mathbf{I}_6$ is the $6 \times 6$ identity
(one-hot matrix).

The 18-move vocabulary consists of 6 faces $\times$ 3 turn types (CW, CCW,
180°).  All move permutations are derived geometrically from face normals and
sticker positions, not hard-coded.  This ensures correctness and makes adding
additional cube sizes straightforward.

Optimal distances are precomputed by BFS from all 24 rotationally equivalent
solved states, producing a table of $88{,}179{,}840$ states mapping
\texttt{state.tobytes()} $\to$ distance.  Seeding from all 24 orientations
removes orientation bias: the solver finds the shortest path to \emph{any}
solved orientation.  BFS takes approximately 25 minutes and is cached to disk.

\subsection{Model Architecture}

We implement \texttt{CubeTransformer} as a \texttt{HookedRootModule} (TransformerLens):

\begin{align*}
\mathbf{h}_0 &= \text{Linear}(144, d_\text{model})(\mathbf{x}) \\
\mathbf{h}_\ell &= \text{TransformerBlock}(\mathbf{h}_{\ell-1}), \quad \ell = 1,\ldots,L \\
\hat{\mathbf{y}} &= \text{Linear}(d_\text{model}, 12)(\text{LayerNorm}(\mathbf{h}_L))
\end{align*}

Default configuration: $d_\text{model}{=}128$, $L{=}4$ blocks, $n_\text{heads}{=}4$,
MLP expansion factor 4 (${\sim}200\text{k}$ parameters).  Since the input is a
single token, attention weights are always 1.0 and computation flows almost
entirely through the MLP sublayers.  All hook points follow TransformerLens
naming (\texttt{hook\_embed}, \texttt{hook\_resid\_mid},
\texttt{hook\_resid\_post}), so \texttt{run\_with\_cache()} works out of the box.

Training uses AdamW ($\text{lr}{=}3{\times}10^{-4}$, $\text{wd}{=}0.01$) with
cosine annealing over 30 epochs, class-weighted cross-entropy to compensate for
the heavy distance-distribution imbalance (distance 1 has 18 states; distance 9
has $\sim$14M).  The dataset contains 50k/5k/5k train/val/test scramble
sequences (${\sim}300\text{k}$ samples each).

\subsection{Linear Probing (Phase 4)}

For each residual stream checkpoint
(\texttt{hook\_embed}, \texttt{resid\_post} layers 0--3), we train:
\begin{itemize}
  \item A \texttt{LogisticRegression} probe per binary feature: 6 face-solved
        bits and 8 corner-orientation bits.
  \item A \texttt{Ridge} regression probe for optimal distance and scramble depth.
\end{itemize}
Probes are fit on a held-out probe split (5k sequences not used in training),
with 5-fold cross-validation for regularization strength.

\subsection{Activation Patching (Phase 5a)}
\label{sec:patch}

We perform two types of patching experiments:

\paragraph{Concept-direction ablation.}
Given a probe weight vector $\mathbf{w}$ for concept $c$, we project it out of
all residual stream activations:
$\hat{\mathbf{h}} = \mathbf{h} - (\mathbf{h} \cdot \hat{\mathbf{w}})\hat{\mathbf{w}}$,
where $\hat{\mathbf{w}} = \mathbf{w}/\|\mathbf{w}\|$.  We then measure whether
model predictions change.  This tests whether the probe direction is
\emph{causally active}.

\paragraph{Counterfactual residual swap.}
For pairs of states from different distance groups ($d_i$ vs.\ $d_j$), we swap
the full residual stream at each layer and measure whether the model's distance
prediction flips.  A 100\% flip rate indicates the residual stream at that layer
contains all the information the model uses to distinguish the two classes.

\subsection{Tuned Lens (Phase 5b)}

For each layer $\ell$, we train an affine probe
$\phi_\ell(\mathbf{h}) = \mathbf{W}_\ell \mathbf{h} + \mathbf{b}_\ell$
to map the intermediate residual stream to the output logit space, then take the
argmax as a layer-$\ell$ prediction.  We compare this to the logit lens (using
the final unembedding $\mathbf{W}_U$ directly on $\mathbf{h}_\ell$) to measure
how much basis rotation is needed to read predictions from early layers.

\subsection{Sparse Autoencoder (Phase 5c)}

For each layer, we train a 2-layer autoencoder with $4{\times}$ expansion (512
features) and L1 penalty on the hidden activations to encourage sparsity:
$\mathcal{L} = \|\mathbf{x} - \hat{\mathbf{x}}\|^2 + \lambda \|\mathbf{z}\|_1$.
Dead features (never active on training data) are resampled from high-loss inputs
every 1,000 steps.  We then compute Pearson correlation between each SAE feature's
activation and known concept labels (face\_solved, optimal\_distance,
corner\_oriented).

\subsection{Text Representation Analysis (Phase 6)}

Motivated by instructor feedback, we evaluate four text representations of the
cube state as inputs to a pre-trained language model (GPT-5 / ChatGPT, accessed
via chat interface).  The representations are:

\begin{enumerate}
  \item \textbf{face\_grid}: An unfolded net diagram showing each face's 2$\times$2
        sticker colors.
  \item \textbf{compact\_string}: A 24-character string, 4 characters per face in
        U/D/F/B/L/R order, reading order within each face.
  \item \textbf{corner\_cubies}: For each corner position, lists which color faces
        which direction (U/D, F/B, L/R).
  \item \textbf{piece\_identity}: Names each corner piece by its solved-state colors
        (\texttt{W-G-R} for the piece belonging at UFR) and states its current
        position and which face its W/Y sticker faces.  This makes the permutation
        and orientation explicit.
\end{enumerate}

All solutions are validated by applying the model's proposed moves to the
scrambled state in our simulator and checking \texttt{is\_solved()}.  We also
test a chain-of-thought (CoT) variant of the piece\_identity prompt, where the
model is instructed to reason step-by-step and conclude with
\texttt{Solution: [moves]}.

% ============================================================
\section{Experiments and Results}
% ============================================================

\subsection{Training}

The model reaches 78.5\% validation accuracy on 12-class distance classification
(random baseline: 8.3\%) in 30 epochs.  Per-class accuracy is high for distances
0--4 and falls for 5--10, reflecting both the class imbalance and the inherent
difficulty of near-maximal scrambles.

\subsection{Linear Probing Results}

\begin{table}[h]
\centering
\small
\begin{tabular}{lcccccc}
\toprule
Feature & Embed & L0 & L1 & L2 & L3 \\
\midrule
\texttt{face\_solved} (avg.\ acc.) & 0.98 & 0.98 & 0.97 & 0.97 & 0.98 \\
\texttt{corner\_oriented} (avg.\ acc.) & 1.00 & 0.97 & 0.96 & 0.94 & 0.92 \\
\texttt{optimal\_distance} (MAE) & 0.87 & 0.40 & 0.39 & 0.39 & 0.38 \\
\bottomrule
\end{tabular}
\caption{Probe accuracy / MAE across residual stream checkpoints.}
\label{tab:probing}
\end{table}

Three findings stand out.  First, \texttt{face\_solved} is 98\% decodable from
\emph{every} layer including the raw embedding---the linear projection from
144-dim one-hot space encodes this feature completely.  Second,
\texttt{corner\_oriented} is perfectly recoverable at the embedding (1.00) but
degrades monotonically to 0.92 at L3, suggesting the model actively transforms
this feature away as it focuses on distance.  Third, the largest improvement in
distance prediction occurs in the very first block (0.87 $\to$ 0.40 MAE), with
minimal improvement in subsequent layers.

\subsection{Activation Patching Results}

\paragraph{Concept-direction ablation.}
Projecting out the \texttt{face\_solved} probe direction at any layer produces
\emph{zero} change in model predictions; the same holds for
\texttt{corner\_oriented}.  Both features are epiphenomenal---encoded but not
causally used.

\paragraph{Counterfactual residual swap.}
Swapping the full residual stream at any layer between distance groups (e.g.,
states with optimal distance 3 vs.\ states with distance 7) produces a 100\%
flip rate, at \emph{every} layer including the raw embedding.  This means the
linear embedding already encodes optimal distance information that is sufficient
to drive the final prediction.

\paragraph{Implication.}
These two results together suggest that the model's linear embedding encodes a
near-complete representation of optimal distance.  The transformer blocks refine
this representation for hard inputs (explaining the tuned-lens results below) but
do not fundamentally restructure it.  The probe-readable features
(\texttt{face\_solved}, \texttt{corner\_oriented}) are by-products of the
representation rather than causes of the output.

\subsection{Tuned Lens Results}

\begin{table}[h]
\centering
\small
\begin{tabular}{lcccccc}
\toprule
Lens & Embed & L0 & L1 & L2 & L3 \\
\midrule
Logit lens  & 0.20 & 0.35 & 0.55 & 0.70 & 0.78 \\
Tuned lens  & 0.60 & 0.70 & 0.75 & 0.79 & 0.81 \\
\bottomrule
\end{tabular}
\caption{Classification accuracy from logit lens and tuned lens at each layer.}
\label{tab:lens}
\end{table}

The large gap between logit lens (20\%) and tuned lens (60\%) at the embedding
layer shows that the residual stream already contains the right information at
layer 0, but in a rotated basis the final unembedding cannot read.  The gap
closes by L3 (78.5\% vs.\ 81.3\%), confirming that the final block aligns the
representation to the unembedding's reading direction.  Distances 1--2 show a
striking phase transition at L2: near-zero accuracy before, $\sim$100\% after.
Distances 5--9 remain hard at every layer, consistent with class imbalance in
training.

This result \emph{reconciles} the patching findings: the blocks are doing
meaningful computation (logit-lens accuracy grows substantially), but they are
rotating the representation rather than adding new information.  The residual
swap's 100\% flip rate reflects that the \emph{content} is present from the
start, while the lens gap reflects that it is not yet in the right format.

\subsection{Sparse Autoencoder Results}

Face and distance features are strongly monosemantic: the top SAE features at
every layer correlate with \texttt{face\_solved} ($r \approx 0.77$--$0.87$) and
\texttt{optimal\_distance} ($r \approx 0.81$--$0.89$).  Symmetric face pairs
(U/D, F/B, L/R) consistently share the same top SAE features, indicating the
model treats geometrically equivalent face pairs identically.  Corner-orientation
features are weaker ($r \approx 0.50$--$0.63$) and strongest at the embedding
layer, consistent with the probing result that corner orientation is encoded early
and then transformed away.  The final layer (L3) is naturally sparser (mean
activation count $\approx 200$ vs.\ $\approx 400$ in earlier layers), consistent
with compression toward a decision.

\subsection{Text Representation Results}

\begin{table}[h]
\centering
\small
\begin{tabular}{lccccc}
\toprule
Representation & $d{=}3$ & $d{=}5$ & $d{=}7$ & Total \\
\midrule
face\_grid         & \checkmark & $\times$ & $\times$ & 1/3 \\
compact\_string    & \checkmark & $\times$ & \checkmark & 2/3 \\
corner\_cubies     & \checkmark & $\times$ & $\times$ & 1/3 \\
\textbf{piece\_identity}  & \checkmark & \checkmark & $\times$ & \textbf{2/3} \\
move\_sequence     & \checkmark & \checkmark & \checkmark & 3/3 \\
piece\_identity+CoT & \checkmark & $\times$ & $\times$ & 1/3 \\
\bottomrule
\end{tabular}
\caption{Solve rate by representation and optimal distance. \texttt{move\_sequence}
is a degenerate baseline (trivial inversion); all others require state understanding.
Tests conducted on GPT-5 (ChatGPT) via chat interface.}
\label{tab:reps}
\end{table}

\paragraph{piece\_identity outperforms position-based representations.}
It is the only state-based representation to consistently solve $d{=}5$
scrambles.  The key difference from \texttt{corner\_cubies} is that it makes the
piece permutation explicit: rather than saying ``the UFR corner has colors
(U:G, F:R, R:W),'' it says ``the W-G-R piece, which belongs at UFR, is currently
at DFL with W facing L.''  This gives the model direct access to the displacement
structure.

\paragraph{$d{=}7$ is a hard wall.}
No state-based representation solved the $d{=}7$ test case reliably.  This
suggests that $d{=}7$ requires tracking how multiple moves propagate through all
eight corners simultaneously---a form of symbolic simulation that the model cannot
do purely from text, regardless of representation quality.

\paragraph{Chain-of-thought did not help.}
The CoT transcript shows the model producing plausible-sounding reasoning
(``I verified this sequence against the exact corner state you gave'') without
actually simulating moves---the solutions it produced were wrong.  The reasoning
is \emph{confabulatory}: narrative about a computation that is not being
performed.  CoT with piece\_identity performed \emph{worse} than direct
piece\_identity (1/3 vs.\ 2/3).

% ============================================================
\section{Discussion: How Results Changed Our Thinking}
% ============================================================

The project began with the hypothesis that the transformer would learn a
hierarchical representation---earlier layers encoding low-level features
(individual sticker colors), later layers encoding higher-level structure
(face solved status, distance to goal).  The results tell a more nuanced story.

The linear embedding already encodes optimal distance nearly completely.  This is
not because the task is easy---78.5\% accuracy on a 12-class problem is non-trivial
---but because the one-hot encoding of a cube state implicitly contains all
distance information in a linearly separable form.  The 144-dim one-hot vector is
highly redundant (physical cube states satisfy hard constraints, so most
144-dim vectors are invalid), and the BFS distance is a smooth function of the
sticker configuration.  The transformer learns to read this function linearly
from the very first layer.

The transformer blocks then \emph{rotate} the representation to align it with the
unembedding direction, as evidenced by the logit/tuned-lens gap.  They do not add
new information but do meaningful work refining the basis.  This is consistent with
Variengien et al.'s prediction that belief states are linearly represented in
residual streams \cite{variengien2024beliefstate}---our model is a direct
instantiation of this, since the cube's belief state for a single-token input is
a point mass.

The text representation experiments added a complementary finding: the
\emph{type} of information matters as much as its presence.  Position-based
representations (face\_grid, corner\_cubies) include all the information needed
to infer the permutation, but the model must compute that inference itself.
piece\_identity pre-computes it.  The performance difference suggests that
pre-trained language models can exploit explicit permutation structure even when
they cannot simulate moves reliably.

% ============================================================
\section{Remaining Work}
% ============================================================

\begin{enumerate}

  \item \textbf{Automated text representation evaluation.}  The current tests use
        only 3 test cases per representation (one per distance level), which is
        insufficient for statistical conclusions.  We plan to run an automated
        evaluation using the Anthropic API (Claude Opus 4.7, as recommended by the
        instructor) with $\geq$10 test cases per distance level, covering $d = 1$
        through $d = 11$.  This will produce a proper solve-rate curve as a
        function of optimal distance.

  \item \textbf{Improved representations.}  We will explore whether augmenting
        piece\_identity with explicit information about which faces each move
        affects (e.g., annotating which pieces are in the U layer) allows the
        model to mentally simulate single moves more reliably.  A hybrid
        representation that provides both the state and a partial BFS hint (e.g.,
        ``these two pieces are the hardest to place and require an anti-Sune
        algorithm'') may substantially improve $d{=}7+$ performance.

  \item \textbf{Mechanistic interpretability of the pretrained model.}  If we
        find a representation where GPT or Claude can reliably solve $d{\leq}7$
        scrambles, the next step is to probe the model's residual stream while it
        processes piece\_identity inputs.  Following Variengien et al.\
        \cite{variengien2024beliefstate}, we predict that the model's belief state
        over cube configurations is linearly encoded---and that the piece\_identity
        representation, by making permutation structure explicit, produces a more
        cleanly linear belief state than position-based representations.

  \item \textbf{Cross-validation of probing results.}  The current probe results
        use a single train/val split.  We will add $k$-fold cross-validation and
        report confidence intervals to strengthen the claims in
        Table~\ref{tab:probing}.

  \item \textbf{More ablation conditions.}  The counterfactual patching experiment
        currently swaps the \emph{full} residual stream.  We will localize this to
        individual components (embedding, each MLP, attention output) to identify
        which architectural component is most causally responsible for the distance
        representation.

\end{enumerate}

% ============================================================
\section{Roadblocks}
% ============================================================

\paragraph{Venv corruption.}
Bumping the package version from 0.1.0 to 0.2.0 without deleting the old
\texttt{.venv} first caused a missing RECORD file warning and an incomplete
install.  Resolved by deleting and recreating the environment.

\paragraph{BFS orientation bias.}
The initial BFS was seeded from a single solved state, so the solver treated the
23 other valid orientations as non-solved.  Solved by seeding from all 24
rotationally equivalent solved states; the regenerated distance table confirms
24 states at distance 0.

\paragraph{Installed-wheel path resolution.}
When installed as a wheel, \texttt{\_\_file\_\_} resolves to
\texttt{site-packages}, not the project root, causing the visualizer's solver
panel to fail to locate \texttt{data/distances\_cache.pkl}.  Fixed by adding a
\texttt{\_find\_project\_root()} function that walks up from \texttt{os.getcwd()}
looking for \texttt{pyproject.toml}.

\paragraph{API access for automated testing.}
Testing text representations against Claude Opus 4.7 (as recommended) requires
an Anthropic API key.  We conducted manual tests via the ChatGPT interface as a
proxy, but the results (3 test cases per representation) have insufficient
statistical power.  Automated testing with more cases is planned.

% ============================================================
\section{Questions for Instructor}
% ============================================================

\begin{enumerate}

  \item \textbf{The confabulatory CoT finding.}  The chain-of-thought transcripts
        show the model claiming to verify its solutions without actually doing
        so.  Is this a finding worth reporting as a standalone result (``pre-trained
        LLMs confabulate cube reasoning under CoT''), or is it too well-known a
        phenomenon to be worth space in the paper?

  \item \textbf{Connecting to belief state geometry.}  Variengien et al.\
        \cite{variengien2024beliefstate} prove that optimal predictors encode
        belief states linearly.  Our from-scratch model is an optimal predictor
        for a task with known ground truth---so in principle the theory
        applies directly.  Is this connection strong enough to justify making it
        a central framing, or is it better treated as motivation rather than
        a formal claim?

  \item \textbf{Statistical power.}  Given only 3 test cases per distance level
        for the text representation experiments, is any conclusion justified other
        than ``piece\_identity is worth further testing''?  Or do the results
        already warrant stronger claims?

  \item \textbf{The 100\% residual-swap flip rate.}  This result (swapping the
        full residual stream at the embedding layer causes 100\% distance-class
        flip) is striking but may be trivially explained by the fact that the
        one-hot encoding fully determines the distance.  Is it still worth
        reporting, or does the trivial explanation make it uninformative?

\end{enumerate}

% ============================================================
\section*{Contributions}
% ============================================================

All code, experiments, and writing were produced by Cole McGuire.  Claude Code
(Anthropic) was used as a coding assistant for implementation and debugging
throughout the project.

% ============================================================
\bibliographystyle{IEEEtran}
\bibliography{references}
\end{document}
